
\documentclass{article}
\usepackage{colortbl}
\usepackage{makecell}
\usepackage{multirow}
\usepackage{supertabular}

\begin{document}

\newcounter{utterance}

\centering \large Interaction Transcript for game `hot\_air\_balloon', experiment `air\_balloon\_survival\_uk\_negotiation\_easy', episode 4 with claude{-}opus{-}4{-}8{-}azure.
\vspace{24pt}

{ \footnotesize  \setcounter{utterance}{1}
\setlength{\tabcolsep}{0pt}
\begin{supertabular}{c@{$\;$}|p{.15\linewidth}@{}p{.15\linewidth}p{.15\linewidth}p{.15\linewidth}p{.15\linewidth}p{.15\linewidth}}
    \# & $\;$A & \multicolumn{4}{c}{Game Master} & $\;\:$B\\
    \hline

    \theutterance \stepcounter{utterance}  
    & & \multicolumn{4}{p{0.6\linewidth}}{
        \cellcolor[rgb]{0.9,0.9,0.9}{
            \makecell[{{p{\linewidth}}}]{
                \texttt{\tiny{[P1$\langle$GM]}}
                \texttt{Ви берете участь у спільній переговорній грі.} \\
\\ 
\texttt{Разом з іншим учасником ви маєте погодити один набір предметів, які буде збережено. Кожен із вас має власне уявлення про те, наскільки кожен предмет важливий для вас (важливість). Ви не знаєте, як інший учасник оцінює предмети. Додатково вам надано зусилля, яких вимагає кожен предмет.} \\
\texttt{Ви можете погодитися на набір лише якщо сумарні зусилля вибраних предметів не перевищують спільний ліміт:} \\
\\ 
\texttt{LIMIT = 3145} \\
\\ 
\texttt{Ось індивідуальні значення зусиль для предметів:} \\
\\ 
\texttt{Item effort = \{"C08": 948, "C55": 118, "A43": 469, "A67": 919, "B42": 284, "A62": 111, "B19": 806, "C43": 47, "C73": 848, "C10": 303, "B36": 13, "A18": 629, "C61": 687, "C67": 15, "C42": 94\}} \\
\\ 
\texttt{Ось ваше особисте уявлення про важливість кожного предмета:} \\
\\ 
\texttt{Item importance values = \{"C08": 138, "C55": 583, "A43": 868, "A67": 822, "B42": 783, "A62": 65, "B19": 262, "C43": 121, "C73": 508, "C10": 780, "B36": 461, "A18": 484, "C61": 668, "C67": 389, "C42": 808\}} \\
\\ 
\texttt{Мета:} \\
\\ 
\texttt{Ваша мета — домовитися про спільний набір предметів, який приносить вам якомога більше користі (тобто максимізує сумарну важливість для ВАС), водночас залишаючись у межах LIMIT. Ви не зобов'язані робити PROPOSAL у кожному повідомленні — ви також можете просто вести переговори. Дозволені всі тактики!} \\
\\ 
\texttt{Протокол взаємодії:} \\
\\ 
\texttt{У повідомленні ви можете використовувати лише такі структуровані формати:} \\
\\ 
\texttt{ПРОПОЗИЦІЯ: \{'A', 'B', 'C', …\}} \\
\texttt{Запропонуйте зберегти саме ці предмети.} \\
\texttt{ВІДМОВА: \{'A', 'B', 'C', …\}} \\
\texttt{Явно відхиліть пропозицію опонента.} \\
\texttt{АРГУМЕНТ: \{'...'\}} \\
\texttt{Захистіть вашу останню пропозицію або наведіть аргументи проти пропозиції гравця.} \\
\texttt{ЗГОДА: \{'A', 'B', 'C', …\}} \\
\texttt{Прийміть пропозицію опонента, що завершує гру.} \\
\\ 
\\ 
\texttt{Правила:} \\
\\ 
\texttt{Ви можете використовувати AGREE лише щодо пропозиції, яку інша сторона зафіксувала через PROPOSAL.} \\
\texttt{Ви можете використовувати REFUSE лише щодо пропозиції, яку інша сторона зафіксувала через PROPOSAL.} \\
\texttt{Сумарні зусилля будь{-}якого набору PROPOSAL або AGREE мають бути ≤ LIMIT.} \\
\texttt{НЕ розкривайте свої приховані оцінки важливості.} \\
\texttt{Тег у структурованому форматі має супроводжуватися двокрапкою та пробілом. Аргумент має бути python{-}множиною, що містить 0 або більше рядків.} \\
\texttt{Отже, це має бути у формі TAG: \{...\}} \\
\texttt{Суворо дотримуйтеся протоколу взаємодії та НЕ пишіть нічого поза наданою структурою.} \\
\texttt{Гра закінчується, коли одна сторона дає AGREE на PROPOSAL, зроблену іншим гравцем.} \\
\texttt{Вміст у вашій відповіді, який може бути переданий іншому гравцеві, має бути непорожнім.} \\
\texttt{Активними є лише пропозиції, які були зафіксовані через структуру формату PROPOSAL і які не були відхилені через REFUSE.} \\
\\ 
            }
        }
    }
    & \\ \\

    \theutterance \stepcounter{utterance}  
    & \multicolumn{4}{p{0.6\linewidth}}{
        \cellcolor[rgb]{0.9,0.9,0.9}{
            \makecell[{{p{\linewidth}}}]{
                \texttt{\tiny{[P1$\rangle$GM]}}
                \texttt{ПРОПОЗИЦІЯ: \{'B36', 'C67', 'C42', 'C55', 'C10', 'B42', 'A43', 'C43', 'C61', 'A67', 'A62'\}} \\
            }
        }
    }
    & & \\ \\

    \theutterance \stepcounter{utterance}  
    & & \multicolumn{4}{p{0.6\linewidth}}{
        \cellcolor[rgb]{0.9,0.9,0.9}{
            \makecell[{{p{\linewidth}}}]{
                \texttt{\tiny{[GM$\rangle$P2]}}
                \texttt{Ви берете участь у спільній переговорній грі.} \\
\\ 
\texttt{Разом з іншим учасником ви маєте погодити один набір предметів, які буде збережено. Кожен із вас має власне уявлення про те, наскільки кожен предмет важливий для вас (важливість). Ви не знаєте, як інший учасник оцінює предмети. Крім того, вам задано зусилля, яких потребує кожен предмет.} \\
\texttt{Ви можете погодити набір лише тоді, коли сумарні зусилля вибраних предметів не перевищують спільний ліміт:} \\
\\ 
\texttt{LIMIT = 3145} \\
\\ 
\texttt{Ось індивідуальні значення зусиль для предметів:} \\
\\ 
\texttt{Item effort = \{"C08": 948, "C55": 118, "A43": 469, "A67": 919, "B42": 284, "A62": 111, "B19": 806, "C43": 47, "C73": 848, "C10": 303, "B36": 13, "A18": 629, "C61": 687, "C67": 15, "C42": 94\}} \\
\\ 
\texttt{Ось ваше особисте уявлення про важливість кожного предмета:} \\
\\ 
\texttt{Item importance values = \{"C08": 138, "C55": 583, "A43": 868, "A67": 822, "B42": 783, "A62": 65, "B19": 262, "C43": 121, "C73": 508, "C10": 780, "B36": 461, "A18": 484, "C61": 668, "C67": 389, "C42": 808\}} \\
\\ 
\texttt{Мета:} \\
\\ 
\texttt{Ваша мета — домовитися про спільний набір предметів, який приносить вам якомога більше користі (тобто максимізує сумарну важливість для ВАС), водночас залишаючись у межах LIMIT. Ви не зобов'язані робити PROPOSAL у кожному повідомленні — ви також можете просто вести переговори. Дозволені всі тактики!} \\
\\ 
\texttt{Протокол взаємодії:} \\
\\ 
\texttt{У повідомленні ви можете використовувати лише такі структуровані формати:} \\
\\ 
\texttt{ПРОПОЗИЦІЯ: \{'A', 'B', 'C', …\}} \\
\texttt{Запропонуйте зберегти рівно ці предмети.} \\
\texttt{ВІДМОВА: \{'A', 'B', 'C', …\}} \\
\texttt{Явно відхиліть пропозицію опонента.} \\
\texttt{АРГУМЕНТ: \{'...'\}} \\
\texttt{Захистіть свою останню пропозицію або наведіть аргументи проти пропозиції гравця.} \\
\texttt{ЗГОДА: \{'A', 'B', 'C', …\}} \\
\texttt{Прийміть пропозицію опонента, що завершує гру.} \\
\\ 
\\ 
\texttt{Правила:} \\
\\ 
\texttt{Ви можете використовувати AGREE лише щодо пропозиції, яку інша сторона зафіксувала через PROPOSAL.} \\
\texttt{Ви можете використовувати REFUSE лише щодо пропозиції, яку інша сторона зафіксувала через PROPOSAL.} \\
\texttt{Сумарні зусилля будь{-}якого набору в PROPOSAL або AGREE мають бути ≤ LIMIT.} \\
\texttt{НЕ розкривайте свої приховані оцінки важливості.} \\
\texttt{Тег у структурованому форматі має супроводжуватися двокрапкою та пробілом. Аргумент має бути python{-}множиною, що містить 0 або більше рядків.} \\
\texttt{Отже, це має бути у формі TAG: \{...\}} \\
\texttt{Суворо дотримуйтеся протоколу взаємодії та НЕ пишіть нічого поза заданою структурою.} \\
\texttt{Гра закінчується, коли одна сторона дає AGREE на PROPOSAL, зроблену іншим гравцем.} \\
\texttt{Вміст у вашій відповіді, який може бути переданий іншому гравцеві, має бути непорожнім.} \\
\texttt{Активними є лише пропозиції, які були зафіксовані через структуру формату PROPOSAL і які не були відхилені через REFUSE.} \\
\\ 
\\ 
\\ 
\texttt{Тепер ви отримаєте перше повідомлення від іншого гравця.} \\
\\ 
\texttt{ПРОПОЗИЦІЯ: \{'B36', 'C67', 'C42', 'C55', 'C10', 'B42', 'A43', 'C43', 'C61', 'A67', 'A62'\}} \\
            }
        }
    }
    & \\ \\

    \theutterance \stepcounter{utterance}  
    & & & \multicolumn{4}{p{0.6\linewidth}}{
        \cellcolor[rgb]{0.9,0.9,0.9}{
            \makecell[{{p{\linewidth}}}]{
                \texttt{\tiny{[GM$\langle$P2]}}
                \texttt{ЗГОДА: \{'B36', 'C67', 'C42', 'C55', 'C10', 'B42', 'A43', 'C43', 'C61', 'A67', 'A62'\}} \\
            }
        }
    }
     \\ \\

    \theutterance \stepcounter{utterance}  
    & & & \multicolumn{2}{p{0.3\linewidth}}{
        \cellcolor[rgb]{0.9,0.9,0.9}{
            \makecell[{{p{\linewidth}}}]{
                \texttt{\tiny{[GM$|$GM]}}
                \texttt{game successful} \\
            }
        }
    }
    & & \\ \\

    \theutterance \stepcounter{utterance}  
    & & & \multicolumn{2}{p{0.3\linewidth}}{
        \cellcolor[rgb]{0.9,0.9,0.9}{
            \makecell[{{p{\linewidth}}}]{
                \texttt{\tiny{[GM$|$GM]}}
                \texttt{end game} \\
            }
        }
    }
    & & \\ \\

\end{supertabular}
}

\end{document}
